\documentclass[11pt]{article}

\usepackage[margin=1in]{geometry}
\usepackage{amsmath,amssymb}
\usepackage{booktabs}
\usepackage{graphicx}
\usepackage{microtype}
\usepackage{caption}
\usepackage[hidelinks]{hyperref}

\newcommand{\stlog}{\mathbf{W}^{(r)}}
\newcommand{\xs}{\mathbf{x}}
\newcommand{\us}{\mathbf{u}}
\newcommand{\ys}{\mathbf{y}}

\title{Toward Real-Time Observability-Aware Control:\\
A Diagnostic Study and a Reformulated Objective}
\author{H.~S.~Helson Go}
\date{\today}

\begin{document}
\maketitle

\begin{abstract}
We study the open problem of solving the observability-aware control (OAC) problem fast
enough for real-time, receding-horizon execution, which the parent work currently performs
only offline. Through systematic profiling we establish that, contrary to the prevailing
assumption, the computational bottleneck is \emph{not} the evaluation of the Short-Time
Local Observability Gramian (STLOG): the entire observability computation accounts for
${\sim}6\%$ of solve time, while the general-purpose nonlinear solver accounts for
${\sim}94\%$. We further find that the objective used in the parent work---the sum over the
horizon of the minimum eigenvalue of the per-node Gramian---is \emph{numerically degenerate}:
its value is floored by the short-time scaling $\lambda_{\min}\sim T^{2r_*+1}$, and its
gradient is effectively zero, so no gradient-based solver can make progress. Replacing it
with a smooth log-determinant surrogate restores a strong gradient and drives trajectories to
full observability. Combined with a lightweight solver (SLSQP) and early termination, the
per-solve wall time on the range-only quadrotor problem drops from ${\sim}10$\,s to
${\sim}100$\,ms (a $10$\,Hz guidance rate) while maintaining full observability. A
closed-loop Extended Kalman Filter (EKF) evaluation on a planar leader--follower analogue
confirms that the fast, early-terminated controller reduces follower-position estimation
error by ${\sim}13\times$ (final) and posterior covariance by ${\sim}20\times$ relative to a
non-OAC baseline, and that early termination preserves the full estimation benefit. The
controller reproduces the theoretically expected observability-optimal \emph{orbiting}
behavior.
\end{abstract}

\paragraph{Reproducibility.} All measurements were taken on a 16-core CPU host
(\texttt{JAX\_PLATFORMS=cpu}, float64), reusing the verified STLOG, Lie-derivative,
integrator, and model implementations from the \texttt{observability\_aware\_control}
companion repository. Raw data and scripts are in \texttt{benchmarks/}, \texttt{experiments/},
and \texttt{results/} of the \texttt{rt-oac} repository.

\section{Problem statement}

The parent work designs control inputs that improve the observability of a leader--follower
quadrotor cooperative-localization system from range and relative-attitude measurements. The
relative state is
\[
\xs = [\,\mathbf{r}^\top,\ \mathbf{q}^\top,\ \mathbf{v}^\top\,]^\top \in \mathbb{R}^{10},
\qquad \mathbf{r}\in\mathbb{R}^3,\ \mathbf{q}\in\mathcal{S}^3,\ \mathbf{v}\in\mathbb{R}^3,
\]
with input $\us\in\mathbb{R}^8$ (leader and follower thrust and body rates) and output
\[
\ys = \left[\,\tfrac{1}{2}d^2,\ \mathbf{q}^\top\,\right]^\top \in \mathbb{R}^5,
\qquad d=\lVert\mathbf{r}\rVert .
\]
Observability over a short horizon $T$ is quantified by the $r$-th order STLOG
\begin{equation}
\stlog = \sum_{i,j=0}^{r}\frac{T^{\,i+j+1}}{(i+j+1)\,i!\,j!}
\, D\!\left(L_\mathbf{f}^i \mathbf{h}\right)^{\!\top} \boldsymbol{\Sigma}^{-1}
D\!\left(L_\mathbf{f}^j \mathbf{h}\right),
\label{eq:stlog}
\end{equation}
evaluated at each node of a predicted trajectory. The observability-predictive controller
(OPC) maximizes a metric of $\{\stlog_k\}_{k=0}^{N-1}$ over a receding horizon subject to
input and inter-vehicle-distance constraints. With $r=5$, $T=0.2$\,s, and $N=20$ the OPC
solve is too slow for online thrust/rate generation and is consequently used offline. The
objective of this study is to identify what actually makes the solve slow and ineffective,
and to make it real-time.

\section{Method}

We constructed a focused testbed (\texttt{rt-oac}) that imports the companion's verified
numerics unchanged and adds: (i) a profiling harness that attributes solve time across
components; (ii) a library of smooth surrogate objectives injected through the cost's existing
\texttt{gramian\_metric} hook; (iii) a controller supporting a selectable solver,
warm-starting, and early termination; and (iv) a closed-loop EKF evaluation. All timed
quantities are JIT-compiled and synchronized with \texttt{block\_until\_ready}; reported wall
times are steady-state (post-compilation). A persistent XLA cache amortizes the one-time
${\sim}200$\,s compilation of the order-$5$ graph.

\section{Results}

\subsection{The optimizer, not the observability computation, is the bottleneck}

Component micro-benchmarks (Table~\ref{tab:components}) show the observability machinery is
cheap once compiled. A single STLOG evaluation costs $52\,\mu$s (versus ${\sim}441$\,ms
\emph{un-compiled}---a ${\sim}6500\times$ penalty that explains the folklore that ``the STLOG
is expensive''). A full objective evaluation over the $N=20$ horizon is $0.5$\,ms; its
reverse-mode gradient $4$\,ms.

\begin{table}[h]
\centering
\caption{Per-component cost (order $5$, $N=20$, JIT-compiled).}
\label{tab:components}
\begin{tabular}{lr}
\toprule
component & cost \\
\midrule
STLOG evaluation (one node) & $52\,\mu$s \\
objective (rollout $+\ 20\times$ STLOG $+$ metric) & $0.5$\,ms \\
gradient (reverse mode) & $4.0$\,ms \\
distance constraint & $0.035$\,ms \\
constraint Jacobian & $0.31$\,ms \\
\bottomrule
\end{tabular}
\end{table}

Attributing one full \texttt{trust-constr} solve ($40$ iterations) shows
$0.68$\,s ($6\%$) inside the JAX callables and $10.6$\,s ($94\%$) inside the solver's own
Python and linear-algebra internals. We verified there is no recompilation across solves and
no measurable NumPy/JAX boundary penalty (the controller's methods cost the same on NumPy or
JAX inputs). The implication is that accelerating the STLOG cannot materially help; the
solver must change.

\subsection{The minimum-eigenvalue objective is numerically degenerate}

At the operating points of interest the per-node Gramian is numerically singular: the minimum
eigenvalue of the sliced (position, velocity) $6\times6$ Gramian is ${\sim}10^{-13}$ across
all horizon nodes, and the same holds for the full tangent-space $9\times9$ Gramian (with the
quaternion gauge removed). Consequently the parent objective $\sum_k\lambda_{\min}(\stlog_k)$
has an automatic-differentiation gradient norm of ${\sim}10^{-13}$, i.e.\ effectively zero.
The cause is intrinsic: the worst-observed direction enters only through high-order Lie
derivatives, whose STLOG weight in~\eqref{eq:stlog} scales as $T^{2r_*+1}=T^{11}$ at order
$5$, i.e.\ ${\sim}2\times10^{-8}$---below the noise floor.

This degeneracy, not solver speed alone, is why the solve is ineffective. Given a zero
gradient, \texttt{trust-constr} exhausts its $40$-iteration budget without improving the
objective, while \texttt{SLSQP} correctly detects the flat landscape and terminates in one
iteration; both end at the worst objective value. This also reconciles an apparent
${\sim}25\times$ ``regression'' ($0.4$\,s $\to 10$\,s): the slow case is the controller
thrashing on the degenerate landscape of a \emph{manifestly unobservable configuration}
(MUC)---here co-hovering / straight co-motion---which the parent thesis enumerates explicitly.
The MUC is expected behavior, not a fault. (A separate latent defect was found in an
in-progress manifold port: a $9\times9$ tangent Gramian sliced with ambient indices clamps an
out-of-bounds index, fabricating a singular Gramian; this has been flagged for repair.)

\subsection{A smooth surrogate restores the gradient and full observability}

We replace the metric with a per-node negative log-determinant,
\begin{equation}
J = -\sum_k \sum_i \log\!\big(\lambda_i(\stlog_k)+\varepsilon\big).
\label{eq:logdet}
\end{equation}
Because $\partial\log(\lambda+\varepsilon)/\partial\lambda=(\lambda+\varepsilon)^{-1}$ is
\emph{largest} exactly where $\lambda\to0$, the log-determinant actively lifts the weak
directions that the minimum-eigenvalue objective cannot. At a non-MUC configuration
(tangential relative velocity, non-identity relative attitude), measuring observability by the
number of well-conditioned directions of the \emph{accumulated} Gramian $\sum_k\stlog_k$
(Table~\ref{tab:objectives}), the log-determinant converts a flat, intractable problem into
one a lightweight solver drives to full observability.

\begin{table}[h]
\centering
\caption{Objective comparison at a non-MUC configuration (SLSQP).}
\label{tab:objectives}
\begin{tabular}{lccc}
\toprule
objective & gradient norm & observable directions (of $6$) & accumulated $\lambda_{\min}$ \\
\midrule
minimum eigenvalue (parent) & ${\sim}10^{-15}$ & $4 \to 4$ (no change) & ${\sim}0$ \\
\textbf{negative log-det} & $\mathbf{12\text{--}15}$ & $\mathbf{4 \to 6}$ \textbf{(full)} & $\mathbf{1.7\text{--}10.6}$ \\
negative soft-min & $0.4$ & $4 \to 6$ & $2.4\text{--}3.9$ \\
\bottomrule
\end{tabular}
\end{table}

\subsection{Real-time solve via a lightweight solver and early termination}

Replacing \texttt{trust-constr} (${\approx}265$\,ms/iteration on this problem) with
\texttt{SLSQP} (${\approx}18$\,ms/iteration), and observing that iterations beyond ${\sim}6$
only refine eigenvalue magnitudes---not the set of observable directions---we early-terminate.
In a true-state receding-horizon rollout of the order-$5$ quadrotor problem from a non-MUC
state (Table~\ref{tab:realtime}), the $10$\,Hz guidance target ($\le 100$\,ms) is met while
maintaining full observability: a ${\sim}100\times$ improvement over the original degenerate
${\sim}10$\,s solve.

\begin{table}[h]
\centering
\caption{Quadrotor closed-loop solve time versus iteration cap.}
\label{tab:realtime}
\begin{tabular}{lcc}
\toprule
iteration cap & median wall / solve & observable directions \\
\midrule
$80$ & $1.0\text{--}1.5$\,s & $6/6$ \\
$12$ & $0.22$\,s & $6/6$ \\
$\mathbf{6}$ & $\mathbf{0.10}$\,\textbf{s} & $\mathbf{6/6}$ \\
\bottomrule
\end{tabular}
\end{table}

\subsection{Closed-loop estimation accuracy (planar EKF)}

To test the decisive question---whether the fast, early-terminated controller preserves
\emph{estimation quality}, not merely a coarse observability count---we used a planar
leader--follower analogue for which a quaternion-free EKF is available. The leader drives
straight; the follower's position is observable only through the inter-vehicle range. We run a
closed loop in which the EKF carries its own estimate (the controller sees only the estimate;
ground truth evolves independently; error accumulates), with a large initial follower-position
error ($2$\,m standard deviation), averaged over $8$ seeds (Table~\ref{tab:ekf}).

\begin{table}[h]
\centering
\caption{Planar closed-loop EKF over $8$ seeds (follower-position error).}
\label{tab:ekf}
\begin{tabular}{lcccc}
\toprule
condition & RMSE (mean\,$\pm$\,std) & final error & terminal cov. & solve \\
\midrule
no-OAC (drive straight) & $1.53\pm0.87$ & $1.52$ & $1.5\times10^{-1}$ & --- \\
OAC, full (cap $60$) & $0.64\pm0.59$ & $0.12$ & $7.6\times10^{-3}$ & $4$\,ms / $10$\,it \\
\textbf{OAC, fast (cap $6$)} & $\mathbf{0.68\pm0.68}$ & $\mathbf{0.11}$ & $\mathbf{8.0\times10^{-3}}$ & $\mathbf{3}$\,\textbf{ms} / $\mathbf{6}$\,\textbf{it} \\
\bottomrule
\end{tabular}
\end{table}

Without OAC the estimator cannot reduce the unobservable initial error (final $\approx$
initial); with OAC the final error falls ${\sim}13\times$, RMSE ${\sim}2.4\times$, and
posterior covariance ${\sim}20\times$. Crucially, the fast controller (cap $6$) matches the
fully converged one on every metric---early termination is safe with respect to estimation
accuracy.

\subsection{Emergent observability-optimal behavior}

The log-determinant controller reproduces the textbook observability-optimal maneuver
(Figure~\ref{fig:traj}). In a planar rollout the follower sweeps a cumulative relative bearing
of $-1613^\circ$ (${\approx}4.5$ revolutions) around the leader---it \emph{orbits}---pressing
against the minimum allowed separation (distance bottoms at the $0.20$\,m bound) to maximize
bearing rate and hence information, while the non-OAC follower sweeps $0^\circ$. This is
consistent with the parent theory and with the known caveat that pure observability
objectives, absent a tracking term, produce operationally awkward (orbiting) trajectories.

\begin{figure}[h]
\centering
\includegraphics[width=\linewidth]{planar_trajectory.png}
\caption{Planar leader--follower trajectory under observability-aware control.
\textit{Left:} world frame---the leader drives straight (black), the non-OAC follower drives
straight and parallel (grey dashed), and the OAC follower traces cycloidal loops (colored by
time). \textit{Right:} leader-relative frame---the OAC follower clearly circles the leader
($\star$), confirming the orbiting behavior.}
\label{fig:traj}
\end{figure}

\section{Discussion}

The study reframes the path to real-time OAC. The dominant levers are, in order:
(1)~\emph{objective reformulation}---the log-determinant restores a usable gradient that the
parent minimum-eigenvalue objective lacks at the relevant operating regime;
(2)~\emph{a lightweight solver}---most of the original solve time was general-purpose
nonlinear-programming overhead, not problem-specific computation; and
(3)~\emph{early termination}---the observability content of a solution saturates within a
handful of iterations. STLOG acceleration and problem transcription, prominent in the initial
plan, proved largely irrelevant because the observability computation is only ${\sim}6\%$ of
the cost.

\paragraph{Limitations.}
(i)~The real-time quadrotor result is validated against an observable-direction count, a
coarse proxy; the rigorous estimation-accuracy validation was performed on the planar
analogue, not the quaternion quadrotor (an error-state EKF for the latter is future work).
(ii)~Measurements are single-host and CPU-only.
(iii)~The observable-direction count thresholds eigenvalues relative to the largest; absolute
conditioning remains poor (the $T^{11}$ floor), which is precisely why the log-determinant
rather than the minimum eigenvalue is the appropriate target.
(iv)~Operating away from MUCs is assumed; escaping MUCs is out of scope.

\section{Conclusion and future work}

Real-time observability-aware control is achievable for the studied system: a smooth
log-determinant objective, a lightweight early-terminated solver, and operation away from
manifestly unobservable configurations reduce the solve from ${\sim}10$\,s to ${\sim}100$\,ms
($10$\,Hz) while a closed-loop EKF confirms the fast controller retains the full estimation
benefit (${\sim}13\times$ lower final error than no OAC). The headline finding is that the
parent minimum-eigenvalue objective is numerically degenerate at the relevant regime and
should be replaced by the log-determinant for any gradient-based real-time use.

Future work: (1)~an error-state / manifold EKF to validate estimation accuracy on the
quaternion quadrotor; (2)~a JAX-native solver to remove the residual ${\sim}13$\,ms/iteration
of solver overhead (the gradient is only $4$\,ms) and target $\ge 50$\,Hz; (3)~a learned
warm-start to approach single-iteration solves; and (4)~a hybrid tracking-plus-observability
objective to temper the orbiting behavior.

\end{document}
